\chapter{Additional Material}
\label{chap:appendix}

This appendix gathers supporting material referenced from the main chapters: the full functional and non-functional requirements registers, the system architecture diagram, application screenshots, the multi-engine benchmark results in tabular form, the interview guide used in the seven semi-structured calls, and the iteration timeline of the eight named iterations.


\section{Functional Requirements Register (FK-01 -- FK-42)}
\label{app:requirements-functional}

\textit{[Placeholder. The full functional requirements register, with each FK-ID, MoSCoW priority, source theme or interview, and implementation status, is to be inserted here as a table. Summary references in the main thesis: \Cref{sec:requirements}, \Cref{sec:eval-traceability}.]}


\section{Non-Functional Requirements Register (IFK-01 -- IFK-16)}
\label{app:requirements-nonfunctional}

\textit{[Placeholder. The full non-functional requirements register, with each IFK-ID, target value, and verification approach, is to be inserted here as a table. Summary references in the main thesis: \Cref{sec:requirements}.]}


\section{System Architecture}
\label{app:architecture}

\textit{[Placeholder. Architecture diagram showing the browser client, the Next.js + tRPC server, the PostgreSQL database, and the three solver subprocesses (greedy heuristic in TypeScript, OR-Tools CP-SAT in Python, Timefold in JVM), with the shared constraint configuration and the multi-tenant \texttt{companyId} boundary. Summary references in the main thesis: \Cref{sec:tech-choice}, \Cref{sec:system-description}.]}


\section{Application Screenshots}
\label{app:screenshots}

\textit{[Placeholder. Screenshots of RP as deployed, showing the planning timeline at \texttt{/ressursplanlegger}, the deviation viewer at \texttt{/avvik}, the dashboard at \texttt{/oversikt}, the drag-and-drop override flow, the time-quality control toggle (fast suggestion vs best-effort plan), and the per-company soft-constraint weight configuration. Summary references in the main thesis: \Cref{sec:system-description}, \Cref{sec:trust-control}.]}


\section{Multi-Engine Benchmark Results}
\label{app:benchmark}

\textit{[Placeholder. Full benchmark results table, with one row per (engine, dataset) pairing, showing scheduled-assignment count and percentage, hard-constraint score, soft-constraint penalty, total lex-tier score, and wall-clock solve time, plus the v0 vs post-rework comparison for the Timefold engine. The reference run is dated 2026-05-02. Summary references in the main thesis: \Cref{sec:eval-benchmark}, \Cref{sec:algorithm}, \Cref{sec:efficiency}.]}


\section{Interview Guide}
\label{app:interview-guide}

\textit{[Placeholder. The five guide topics used to anchor the seven semi-structured calls (current tools and workflow, assignment criteria, sick-leave handling, attitudes to automation, adoption criteria), with the broad opening question for each topic and the prompts used where coordinator answers left an obvious gap. Summary references in the main thesis: \Cref{sec:data-collection}.]}


\section{Iteration Timeline}
\label{app:iterations}

\textit{[Placeholder. Gantt or table of the eight named iterations (schema and timeline scaffolding, single-engine baseline, constraint generalisation, multi-engine selection and benchmarking, HITL interface, real-time deviation detection, plan-time vs plan-quality tradeoff, configurable soft-constraint weights), with sprint boundaries and the dated capability milestones from the project's working notes. Summary references in the main thesis: \Cref{sec:iterations}, \Cref{sec:sprint-summary}.]}


\section{Project Process Artefacts}
\label{app:process}

\textit{[Placeholder. Supporting process documentation as agreed with the supervisor and Admmit: GitHub milestones for the eight iterations, time-accounting summary by author and task category, decision log highlights, and any meeting notes the authors choose to include. Summary references in the main thesis: \Cref{sec:key-decisions}, \Cref{sec:time-tracking}.]}
